% Keep the existing bold commands on the OpenType mathematical alphabets.
\ExplSyntaxOn
\DeclareRobustCommand{\bm}[1]{%
  \regex_match:nnTF { \A [A-Za-z]+ \Z } {#1}
    { \symbfit{#1} }
    { \symbf{#1} }
}
\ExplSyntaxOff
\DeclareRobustCommand{\boldsymbol}[1]{\bm{#1}}

\numberwithin{equation}{chapter}
\counterwithin{figure}{chapter}
\AtBeginDocument{%
  \theoremstyle{plain}%
  \theoremheaderfont{\sffamily\bfseries\thmheiti}%
  \theorembodyfont{\fangsong}%
  \theoremseparator{\enspace}%
  \theoremsymbol{}%
  \newtheorem{RiemannDefinition}{Definition}[section]%
  \renewcommand{\theRiemannDefinition}{\arabic{section}.\arabic{RiemannDefinition}}%
  \newtheorem{RiemannTheorem}[RiemannDefinition]{Theorem}%
  \newtheorem{RiemannLemma}[RiemannDefinition]{Lemma}%
  \newtheorem{RiemannProposition}[RiemannDefinition]{Proposition}%
  \newtheorem{RiemannCorollary}[RiemannDefinition]{Corollary}%
  \theorembodyfont{\normalfont}%
  \newtheorem{RiemannExample}[RiemannDefinition]{Example}%
  \newtheorem{RiemannRemark}[RiemannDefinition]{Remark}%
  \newtheorem{RiemannExercise}[RiemannDefinition]{Exercise}%
  \theoremstyle{nonumberplain}%
  \theoremheaderfont{\sffamily\bfseries\CJKfamily{pffont}}%
  \theoremsymbol{\openbox}%
  \newtheorem{RiemannProof}{Proof}%
  \setlength{\abovedisplayskip}{8pt}%
  \setlength{\belowdisplayskip}{8pt}%
  \setlength{\abovedisplayshortskip}{8pt}%
  \setlength{\belowdisplayshortskip}{8pt}%
}
\makeatletter
% For stages of the translation where the target is not yet in the manuscript,
% retain a verified link to its page in the original. No provisional aux labels.
\newcommand{\RGRegisterSourceRef}[3]{%
  \expandafter\def\csname RGsourcepage:#1\endcsname{#3}%
  \expandafter\def\csname RGsourcevalue:#1\endcsname{#2}%
}
\newcommand{\RGref}[1]{%
  \@ifundefined{r@#1}{%
    \@ifundefined{RGsourcepage:#1}{\ref{#1}}{%
      \href[page=\csname RGsourcepage:#1\endcsname]{../../input/source/An Introduction to Riemannian Geometry.pdf}{\csname RGsourcevalue:#1\endcsname}%
    }%
  }{\ref{#1}}%
}
\newcommand{\RGbiblabel}[2][]{%
  \phantomsection
  \begingroup\def\@currentlabel{#2}\label{bib:#1#2}\endgroup
}
\newcommand{\RGsolution}[2]{%
  \item[(\ref{#1})]\leavevmode\phantomsection
  \begingroup\edef\@currentlabel{\getrefnumber{#1}}\label{sol:#2}\endgroup
}
\makeatother
\newcommand{\RGcite}[1]{[\RGref{bib:#1}]}
\newcommand{\RGReferenceSection}{%
  \section*{参考文献}%
  \phantomsection
  \addcontentsline{toc}{section}{参考文献}%
}
\newcommand{\RGfigure}[3][.6\textwidth]{%
  \begin{figure}[htbp]
  \centering
  \includegraphics[width=#1]{figures/#2.pdf}
  \caption{#3}\label{fig:#2}
  \end{figure}%
}
